% GENERATED FILE. Edit canonical lessons, then run npm run generate:books.
% canonical-source-hash: fnv1a64:60a58e7c1ecb1235
% canonical-lessons: MW-C249-holi, MW-C249-gulal, MW-C249-sugan, MW-C249-divali, MW-C249-puja

\chapter{Holi and Diwali}
\label{ch:mw-holi-and-diwali}
\hlchaptermodality{249}

\begin{chapteropening}
\textbf{By the end of this chapter:} \emph{I can say which festival comes when.}

Complete the last lesson of chapter 249: I can say which festival comes when.
\end{chapteropening}

\section[\texorpdfstring{\mw{होळी} (hoḷī)}{होळी (hoḷī)}]{\mw{होळी} (hoḷī) — the spring festival of colours, Holi}

\label{lesson:MW-C249-holi}



\begin{quote}
Before the new one: say the Marwadi for to dress up, then the Marwadi for a bridegroom's garland.
\end{quote}

\subsection*{You'll want to know: \mw{होळी}}
\textbf{\mw{होळी}} — \emph{hoḷī} — \textquotedblleft{}the spring festival of colours, Holi\textquotedblright{}.

\begin{quote}
\mw{होळी} \mw{बसंत} \mw{मांय} \mw{आवै} \mw{है।}
\end{quote}

\begin{quote}
\emph{hoḷī basant mā̃y āvai hai.}
\end{quote}

\begin{quote}
Holi comes in spring.
\end{quote}

A feminine noun.

\begin{cousinweb}
\textbf{\mw{होळी}} (\emph{hoḷī}) is inherited from Sanskrit \emph{holā}, \textquotedblleft{}the spring festival\textquotedblright{}. CDIAL gives no Marwari form in this entry; it lists Hindi \emph{holī} and Gujarati \emph{hoḷī}, with this book's \mw{ळ} (CDIAL 14182).
\end{cousinweb}

\begin{grammarlens}[title={What you've built}]
One more word for saying what happened.
\end{grammarlens}

\subsection*{Your turn}
\begin{itemize}
\raggedright
  \item \emph{Say it:} \emph{hoḷī}
  \item \emph{Say it:} \emph{hoḷī}, once more
  \item \emph{Say it:} \emph{hoḷī basant mā̃y āvai hai.} — Holi comes in spring.
  \item \emph{Recall:} say the Marwadi for to dress up, then the Marwadi for a bridegroom's garland, then say \emph{hoḷī} again
\end{itemize}

\begin{tcolorbox}[breakable,colback=teal!4,colframe=teal!35!black,title={Before you move on}]
What does \mw{होळी} mean? (\textquotedblleft{}The spring festival of colours, Holi\textquotedblright{}.) Say it once more.
\end{tcolorbox}
\section[\texorpdfstring{\mw{गुलाल} (gulāl)}{गुलाल (gulāl)}]{\mw{गुलाल} (gulāl) — coloured powder for Holi}

\label{lesson:MW-C249-gulal}



\begin{quote}
Before the new one: say the Marwadi for a bridegroom's garland, then the Marwadi for Holi.
\end{quote}

\subsection*{You'll want to know: \mw{गुलाल}}
\textbf{\mw{गुलाल}} — \emph{gulāl} — \textquotedblleft{}coloured powder for Holi\textquotedblright{}.

\begin{quote}
\mw{गुलाल} \mw{लाल} \mw{अर} \mw{गुलाबी} \mw{है।}
\end{quote}

\begin{quote}
\emph{gulāl lāl ar gulābī hai.}
\end{quote}

\begin{quote}
The Holi powder is red and pink.
\end{quote}

A masculine noun.

\begin{cousinweb}
\textbf{\mw{गुलाल}} (\emph{gulāl}) is usually traced to Persian \emph{gulāl}, a red powder; no dictionary we can open gives an older Indian form.
\end{cousinweb}

\begin{grammarlens}[title={What you've built}]
One more word for saying what happened.
\end{grammarlens}

\subsection*{Your turn}
\begin{itemize}
\raggedright
  \item \emph{Say it:} \emph{gulāl}
  \item \emph{Say it:} \emph{gulāl}, once more
  \item \emph{Say it:} \emph{gulāl lāl ar gulābī hai.} — The Holi powder is red and pink.
  \item \emph{Recall:} say the Marwadi for a bridegroom's garland, then the Marwadi for Holi, then say \emph{gulāl} again
\end{itemize}

\begin{tcolorbox}[breakable,colback=teal!4,colframe=teal!35!black,title={Before you move on}]
What does \mw{गुलाल} mean? (\textquotedblleft{}Coloured powder for Holi\textquotedblright{}.) Say it once more.
\end{tcolorbox}
\section[\texorpdfstring{\mw{सुगन} (sugan)}{सुगन (sugan)}]{\mw{सुगन} (sugan) — a good omen; an auspicious sign}

\label{lesson:MW-C249-sugan}



\begin{quote}
Before the new one: say the Marwadi for Holi, then the Marwadi for coloured powder for Holi.
\end{quote}

\subsection*{You'll want to know: \mw{सुगन}}
\textbf{\mw{सुगन}} — \emph{sugan} — \textquotedblleft{}a good omen; an auspicious sign\textquotedblright{}.

\begin{quote}
\mw{आज} \mw{घणो} \mw{आछो} \mw{सुगन} \mw{है।}
\end{quote}

\begin{quote}
\emph{āj ghaṇo āchho sugan hai.}
\end{quote}

\begin{quote}
Today there is a very good omen.
\end{quote}

A masculine noun.

\begin{cousinweb}
\textbf{\mw{सुगन}} (\emph{sugan}) is probably a later borrowing of Sanskrit \emph{śakuna}, \textquotedblleft{}a bird\textquotedblright{}, and so \textquotedblleft{}an omen\textquotedblright{} read from birds. CDIAL gives no Marwari form in this entry; it lists Hindi \emph{saun}, \textquotedblleft{}good omen\textquotedblright{}, a more worn-down shape of the same word (CDIAL 12245).
\end{cousinweb}

\begin{grammarlens}[title={What you've built}]
One more word for saying what happened.
\end{grammarlens}

\subsection*{Your turn}
\begin{itemize}
\raggedright
  \item \emph{Say it:} \emph{sugan}
  \item \emph{Say it:} \emph{sugan}, once more
  \item \emph{Say it:} \emph{āj ghaṇo āchho sugan hai.} — Today there is a very good omen.
  \item \emph{Recall:} say the Marwadi for Holi, then the Marwadi for coloured powder for Holi, then say \emph{sugan} again
\end{itemize}

\begin{tcolorbox}[breakable,colback=teal!4,colframe=teal!35!black,title={Before you move on}]
What does \mw{सुगन} mean? (\textquotedblleft{}A good omen; an auspicious sign\textquotedblright{}.) Say it once more.
\end{tcolorbox}
\section[\texorpdfstring{\mw{दिवाळी} (divāḷī)}{दिवाळी (divāḷī)}]{\mw{दिवाळी} (divāḷī) — the festival of lamps, Diwali}

\label{lesson:MW-C249-divali}



\begin{quote}
Before the new one: say the Marwadi for coloured powder for Holi, then the Marwadi for a good omen.
\end{quote}

\subsection*{You'll want to know: \mw{दिवाळी}}
\textbf{\mw{दिवाळी}} — \emph{divāḷī} — \textquotedblleft{}the festival of lamps, Diwali\textquotedblright{}.

\begin{quote}
\mw{दिवाळी} \mw{घणी} \mw{चोखी} \mw{लागै} \mw{है।}
\end{quote}

\begin{quote}
\emph{divāḷī ghaṇī chokhī lāgai hai.}
\end{quote}

\begin{quote}
Diwali feels wonderful.
\end{quote}

A feminine noun.

\begin{cousinweb}
\textbf{\mw{दिवाळी}} (\emph{divāḷī}) is inherited from Sanskrit \emph{dīpāvali}, \textquotedblleft{}a row of lamps\textquotedblright{}. CDIAL gives no Marwari form in this entry; it lists Gujarati \emph{divāḷī} (CDIAL 6358).
\end{cousinweb}

\begin{grammarlens}[title={What you've built}]
One more word for saying what happened.
\end{grammarlens}

\subsection*{Your turn}
\begin{itemize}
\raggedright
  \item \emph{Say it:} \emph{divāḷī}
  \item \emph{Say it:} \emph{divāḷī}, once more
  \item \emph{Say it:} \emph{divāḷī ghaṇī chokhī lāgai hai.} — Diwali feels wonderful.
  \item \emph{Recall:} say the Marwadi for coloured powder for Holi, then the Marwadi for a good omen, then say \emph{divāḷī} again
\end{itemize}

\begin{tcolorbox}[breakable,colback=teal!4,colframe=teal!35!black,title={Before you move on}]
What does \mw{दिवाळी} mean? (\textquotedblleft{}The festival of lamps, Diwali\textquotedblright{}.) Say it once more.
\end{tcolorbox}
\section[\texorpdfstring{\mw{पूजा} (pūjā)}{पूजा (pūjā)}]{\mw{पूजा} (pūjā) — worship; a prayer ceremony}

\label{lesson:MW-C249-puja}



\begin{quote}
Before the new one: say the Marwadi for a good omen, then the Marwadi for Diwali.
\end{quote}

\subsection*{You'll want to know: \mw{पूजा}}
\textbf{\mw{पूजा}} — \emph{pūjā} — \textquotedblleft{}worship; a prayer ceremony\textquotedblright{}.

\begin{quote}
\mw{मां} \mw{रोज} \mw{पूजा} \mw{करै} \mw{है।}
\end{quote}

\begin{quote}
\emph{mā̃ roj pūjā karai hai.}
\end{quote}

\begin{quote}
Mother does the worship every day.
\end{quote}

A feminine noun.

\begin{cousinweb}
\textbf{\mw{पूजा}} (\emph{pūjā}) is a tatsama, the Sanskrit word \emph{pūjā}, \textquotedblleft{}honour, worship\textquotedblright{}, taken over unchanged (Monier-Williams).
\end{cousinweb}

\begin{grammarlens}[title={What you've built}]
One more word for saying what happened.
\end{grammarlens}

\subsection*{Your turn}
\begin{itemize}
\raggedright
  \item \emph{Say it:} \emph{pūjā}
  \item \emph{Say it:} \emph{pūjā}, once more
  \item \emph{Say it:} \emph{mā̃ roj pūjā karai hai.} — Mother does the worship every day.
  \item \emph{Recall:} say the Marwadi for a good omen, then the Marwadi for Diwali, then say \emph{pūjā} again
\end{itemize}

\begin{tcolorbox}[breakable,colback=teal!4,colframe=teal!35!black,title={Before you move on}]
What does \mw{पूजा} mean? (\textquotedblleft{}Worship; a prayer ceremony\textquotedblright{}.) Say it once more.
\end{tcolorbox}
